\documentclass[10pt, a4paper]{article}
\usepackage[a4paper, top=1.3cm, bottom=1.3cm, left=1.35cm, right=1.35cm,
            headheight=20pt, headsep=7pt, footskip=13pt]{geometry}

\usepackage{amsmath,amssymb}
\usepackage{fontspec}
\IfFileExists{fonts/times.ttf}{%
  \setmainfont{times.ttf}[Path=fonts/, BoldFont=timesbd.ttf,
                          ItalicFont=timesi.ttf, BoldItalicFont=timesbi.ttf]%
}{%
  \setmainfont{TeX Gyre Termes}%
}

\usepackage{xcolor}
\usepackage{booktabs}
\usepackage{tabularx}
\usepackage{colortbl}
\usepackage{fancyhdr}
\usepackage{lastpage}
\usepackage{titlesec}
\usepackage[most]{tcolorbox}
\usepackage[hidelinks]{hyperref}

% IEEE / Nature / Viridis Scientific Palette
\definecolor{navyblue}{HTML}{1B365D}
\definecolor{slate}{HTML}{334155}
\definecolor{softblue}{HTML}{EBF2F7}
\definecolor{softteal}{HTML}{E6F2F2}
\definecolor{softred}{HTML}{F9EBEB}
\definecolor{accentteal}{HTML}{1F7A7A}
\definecolor{accentblue}{HTML}{2B6CB0}
\definecolor{accentred}{HTML}{A82020}
\definecolor{lightgray}{HTML}{F8FAFC}
\definecolor{bordergray}{HTML}{94A3B8}

\titleformat{\section}{\fontsize{11.5pt}{13.5pt}\selectfont\bfseries\color{navyblue}}
  {\thesection.}{0.3em}{}[{\color{accentteal}\titlerule[0.7pt]}]
\titleformat{\subsection}{\fontsize{10pt}{12pt}\selectfont\bfseries\color{accentblue}}{\thesubsection}{0.3em}{}

\titlespacing*{\section}{0pt}{2.5pt}{1.5pt}
\titlespacing*{\subsection}{0pt}{1.5pt}{1pt}

\setlength{\parindent}{0pt}
\setlength{\parskip}{1.2pt}
\linespread{0.97}

\pagestyle{fancy}
\fancyhf{}
\fancyhead[L]{\fontsize{8pt}{10pt}\selectfont\color{navyblue}\textbf{RahiScan Viva Study Guide} \(\cdot\) \textcolor{accentblue}{Muhammad Saad Hussain (SU92-BSCSM-F25-260)}}
\fancyhead[R]{\fontsize{8pt}{10pt}\selectfont\color{slate}Intro to Data Science \(\cdot\) Fall 2026}
\fancyfoot[L]{\fontsize{7.5pt}{9.5pt}\selectfont\color{slate}Superior University Gold Campus \(\cdot\) Group Proposal Defense}
\fancyfoot[R]{\fontsize{8pt}{10pt}\selectfont\color{navyblue}\textbf{Page \thepage\ of \pageref{LastPage}}}
\renewcommand{\headrule}{{\color{accentteal}\hrule height 0.6pt}}
\renewcommand{\footrule}{{\color{bordergray}\hrule height 0.5pt}}

\newtcolorbox{qa}[1][]{enhanced, colback=lightgray, colframe=navyblue, boxrule=0.5pt,
  arc=2.5pt, left=6pt, right=6pt, top=2.5pt, bottom=2.5pt, fonttitle=\bfseries\footnotesize,
  coltitle=white, colbacktitle=navyblue, title={#1}, before skip=1.5pt, after skip=1.5pt}

\newtcolorbox{speaking}[1][]{enhanced, colback=softblue, colframe=accentblue, boxrule=0.5pt,
  arc=2.5pt, left=6pt, right=6pt, top=2.5pt, bottom=2.5pt, fonttitle=\bfseries\footnotesize,
  coltitle=white, colbacktitle=accentblue, title={#1}, before skip=1.5pt, after skip=1.5pt}

\newtcolorbox{overviewbox}[1][]{enhanced, colback=softteal, colframe=accentteal, boxrule=0.5pt,
  arc=2.5pt, left=6pt, right=6pt, top=2.5pt, bottom=2.5pt, fonttitle=\bfseries\footnotesize,
  coltitle=white, colbacktitle=accentteal, title={#1}, before skip=1.5pt, after skip=1.5pt}

\begin{document}

% Top Banner
\begin{tcolorbox}[enhanced, colback=navyblue, colframe=navyblue, arc=3pt, left=8pt, right=8pt, top=4pt, bottom=4pt]
  {\fontsize{15pt}{17pt}\selectfont\color{white}\bfseries RahiScan \(\cdot\) Student Viva Study Guide}\\[1.5pt]
  {\fontsize{10pt}{12pt}\selectfont\color{softteal}\textbf{Member 1: Muhammad Saad Hussain} (Roll No: \texttt{SU92-BSCSM-F25-260})}\\[1pt]
  {\fontsize{8.5pt}{10.5pt}\selectfont\color{white!85!softblue}\textbf{Assigned Role:} Architect \& Product Lead \(\cdot\) \textbf{Primary Stages:} Case Study, Define, Collect, Deploy}
\end{tcolorbox}

\begin{overviewbox}[Quick 30-Second Group Overview (All Members Must Know This)]
\footnotesize
\begin{itemize}\setlength{\itemsep}{0.8pt}\setlength{\parskip}{0pt}
  \item \textbf{What is RahiScan?} A crowd-sourced mobile app proposal for students commuting by \textbf{Qingqi rickshaw} to \textbf{Superior University Gold Campus} (Raiwind Road) and across Lahore corridors.
  \item \textbf{The Core Problem:} Qingqi rickshaws lack fare meters. Curbside passengers face two unknowns: (1) Is the quoted fare fair or inflated? (2) Is the driver normal or habitually slow?
  \item \textbf{How It Works:} Snap one photo before boarding. The app reads the \textbf{number plate} and \textbf{driver's face}, checks past logs, and returns a scorecard in \textbf{\(< 2\) seconds}. During the trip, it logs 5-second GPS telemetry.
  \item \textbf{Status:} Comprehensive research proposal backed by a 4-week simulated pilot (\textbf{1,224 rides}, \textbf{40 drivers}).
\end{itemize}
\end{overviewbox}

\begin{tcolorbox}[enhanced, colback=white, colframe=bordergray, boxrule=0.4pt, arc=2.5pt, left=6pt, right=6pt, top=2.5pt, bottom=2.5pt]
  \footnotesize \textbf{Your Role in Simple Words:} You open the group viva by setting the scene: narrating our morning commute story, explaining how we converted commuter frustration into quantitative metrics (\textbf{Define}), describing how data is collected with biometric privacy safeguards (\textbf{Collect}), and concluding with how the app is safely deployed at the roadside (\textbf{Deploy}).
\end{tcolorbox}

\section{What to Say: Step-by-Step Viva Presentation Script}

\subsection{A. The Case Study Story \& Lifecycle Mapping (Report Page 2)}
\begin{speaking}[Speaking Script: Case Study Story]
\footnotesize
``Our project originated from a real Monday morning commute at 7:40 AM from Johar Town to an 8:30 AM lab at \textbf{Superior University Gold Campus} on Raiwind Road (9.6 km away). The first Qingqi driver quoted \textbf{140 PKR}; we bargained down to \textbf{110 PKR} and boarded. He crawled along the left lane and stopped for tea at a dhaba, making us twelve minutes late. The next morning, a different driver completed the identical trip for the standard \textbf{80 PKR} in twenty minutes despite heavier traffic (similarly, \textbf{Raiwind to Gold Campus} is \textbf{120 PKR} for 16.1 km, and inner-city trips are 40--85 PKR).

From the curb, both Qingqis looked identical. We designed \textbf{RahiScan}: snapping one photo before boarding (Figure 1) fuses plate OCR with facial recognition because vehicles in Lahore are frequently shift-shared. From this real story, we mapped out all seven Data Science Lifecycle stages (Figure 2, Table 1).''
\end{speaking}

\subsection{B. Stage 1 --- Define: Turning Opinions into Testable Targets (Report Page 3)}
\begin{speaking}[Speaking Script: Define Stage]
\footnotesize
``In the \textbf{Define} stage, we recognized that `slow' and `overpriced' are purely subjective without context. A driver going 10 km/h in peak Mall Road gridlock is just stuck in traffic. We formulated two dimensionless ratios centered on \textbf{1.0 as normal}:
\begin{enumerate}\setlength{\itemsep}{0.8pt}\setlength{\parskip}{0pt}
  \item \textbf{Driving Pace Index (\(\Pi\)):} Driver speed divided by crowd-sourced median speed \(v_{\mathrm{base}}\) of peer rickshaws on that corridor and hour. \(\Pi < 0.75\) earns a \textbf{Habitually Slow} badge (Table 2).
  \item \textbf{Fare Fairness Ratio (\(\Phi\)):} Quoted fare divided by predicted fair fare (\(\approx 80\) PKR for 9.6 km; \(\approx 120\) PKR Raiwind to Gold Campus). \(\Phi > 1.25\) earns an \textbf{Overpriced} badge.''
\end{enumerate}
\end{speaking}

\subsection{C. Stage 2 --- Collect: Gathering Telemetry Responsibly (Report Page 3)}
\begin{speaking}[Speaking Script: Collect Stage]
\footnotesize
``In the \textbf{Collect} stage, since no public dataset links Lahore Qingqi drivers to GPS telemetry, we designed a 4-week pilot with 85 student volunteers commuting to \textbf{Superior University Gold Campus}. We simulated \textbf{1,224 rides}, \textbf{1,200 scans}, and \textbf{411,829 GPS pings} across \textbf{40 drivers} and 34 vehicles (Table 3, Figure 4).

We collect 5 streams: curbside photo (plate + face), quoted fare, 5-second GPS + accelerometer pings, corridor fare benchmarks, and shift registries. For privacy, raw photos are processed instantly on-device and discarded; only irreversible 128-dimensional mathematical embeddings are stored in our secure database.''
\end{speaking}

\newpage

\subsection{D. Stage 7 --- Deploy: Curbside Decision \& Operational Safeguards (Report Page 6)}
\begin{speaking}[Speaking Script: Deploy Stage]
\footnotesize
``In the \textbf{Deploy} stage (Figure 9), the passenger uses three clean screens: curbside scan, live ride telemetry, and an instant scorecard returned in \textbf{under two seconds} (alerting students when a 140 PKR quote exceeds the 80 PKR benchmark).

To protect driver livelihoods from false accusations, we engineered four operational safeguards (Table 6):
\begin{enumerate}\setlength{\itemsep}{0.8pt}\setlength{\parskip}{0pt}
  \item \textbf{The 5-Ride Rule:} Badges are suppressed until a driver logs \(\ge 5\) verified trips, absorbing single-trip traffic bad luck.
  \item \textbf{Shift-Shared Protection:} Profiles are indexed by composite \texttt{Plate\#FaceID}, so morning and night drivers sharing one vehicle have distinct records.
  \item \textbf{Recency Weighting \& Dispute Reviews:} Recent rides carry greater weight, and drivers can scan their own Qingqi to review scores and contest disputed trips.
  \item \textbf{Monthly Model Recalibration:} The fare regressor is updated monthly against official rolling petrol/LPG prices.''
\end{enumerate}
\end{speaking}

\section{Probable Teacher Questions \& Winning Answers}

\begin{qa}[Q1: Why do you need both number plate AND driver face? Isn't the plate enough?]
\footnotesize
\textbf{Answer:} In Lahore, Qingqi rickshaws are commonly rented out across morning and evening shifts. In our pilot, 6 out of 34 vehicles were shift-shared. If we tracked plates alone, an honest morning driver charging 80 PKR and an overpriced evening driver charging 140 PKR would be blended together. Facial recognition ensures each individual driver builds their own genuine reputation.
\end{qa}

\begin{qa}[Q2: What if a driver is delayed by an accident or police barricade? Do they get penalized?]
\footnotesize
\textbf{Answer:} No. First, the Pace Index \(\Pi\) normalizes speed against the ambient baseline \(v_{\mathrm{base}}\) of peer rickshaws on that road and hour (everyone slows down together). Second, badges require a \textbf{median across at least 5 verified rides}, meaning one unlucky delay has zero lasting effect on their rating.
\end{qa}

\begin{qa}[Q3: How does the student benefit in real life right at the roadside?]
\footnotesize
\textbf{Answer:} Within two seconds of snapping the photo, the student sees the fair benchmark range (75--85 PKR) and the driver's pace badge. If the driver demands 140 PKR, the student can either negotiate with confidence using the app's benchmark or politely hail the next Qingqi.
\end{qa}

\section{Saad's Key Numbers Cheat-Sheet}
\footnotesize
\begin{tabularx}{\textwidth}{|>{\bfseries\color{navyblue}}p{3.8cm}|X|>{\raggedleft\arraybackslash}p{3.0cm}|}
\hline
\rowcolor{navyblue}\textcolor{white}{\textbf{Metric}} & \textcolor{white}{\textbf{Context}} & \textcolor{white}{\textbf{Value to State}} \\ \hline
Campus \& Corridors & Superior University Gold Campus & 8 Routes, 85 Volunteers \\ \hline
Typical Qingqi Fares & Johar Town / Thokar to Gold Campus (9.6 km) & \textbf{80 PKR} (Fair benchmark) \\ \hline
Raiwind to Gold Campus & Longer morning commute corridor (16.1 km) & \textbf{120 PKR} (Fair benchmark) \\ \hline
Simulated Pilot Volume & 4 Weeks, 40 Drivers, 34 Qingqis & \textbf{1,224} Raw, \textbf{1,071} Clean \\ \hline
Curbside Latency & Time to return scorecard to student & \textbf{\(< 2\) Seconds} \\ \hline
Badge Minimum & Required verified trips before badge displays & \textbf{\(\ge 5\) verified rides} \\ \hline
\end{tabularx}

\end{document}
